\begin{question}{30}{
    Een officier logistiek bij het Commando Materieel en IT (COMMIT) voert een onderzoek uit naar de leveringsbetrouwbaarheid 
    van twee leveranciers, $A$ en $B$, van kritieke onderdelen voor Chinook-transporthelikopters.
    Hieronder staat een dataset weergegeven in een tabel, waarbij voor beide leveranciers de levertijd (in weken) is geregistreerd voor een aantal orders voor kritieke helikopteronderdelen:     
    \begin{center}
        \begin{tabular}{m{2.7cm}|cccccccccccc}
            \toprule
                \textbf{Levertijd $A$} \newline \textbf{(in weken)} & $12.1$ & $11.9$ & $12.3$ & $11.8$ & $12.2$ & $12.0$ & $11.9$ & $12.4$ & $11.7$ & $12.1$ & $12.0$ & $11.8$\\
            \midrule
                \textbf{Levertijd $B$} \newline \textbf{(in weken)} & $11.9$ & $7.3$ & $14.7$ & $6.0$ & $13.6$ & $8.0$ & $9.3$ & $5.0$ & $7.6$ & $11.4$ & & \\
            \bottomrule
        \end{tabular}
    \end{center}
    Het doel van COMMIT is om op basis van de resultaten van dit onderzoek een keuze te maken voor één van beide leveranciers voor het leveren van Chinook-onderdelen.
    Je mag in deze vraag bij hypothesetoetsen gebruik maken van een significantieniveau $\alpha = 0.05$.
}

    \subquestion{7}{
        Bereken de steekproefgemiddelden en de steekproefvarianties van de levertijden van beide leveranciers.
    }
    \solution{
        We berekenen het steekproefgemiddelde $\overline{x_A}$ voor de $n = 12$ levertijden van leverancier $A$ als volgt:
        \begin{align*}
            \overline{x_A} = \frac{a_1+a_2+\ldots+a_n}{n} = \frac{ 12.1 + 11.9 + \ldots + 11.8 }{ 12 } \approx 12.0167.\rubric{1}
        \end{align*}
    
        We berekenen de steekproefvariantie $s_A^2$ als volgt:
        \begin{align*}
            s_A^2 &= \frac{ (a_1 - \overline{x_A})^2 + (a_2 - \overline{x_A})^2 + \ldots + (a_n - \overline{x_A})^2 }{ n - 1 } \\
                &= \frac{ (12.1 - 12.0167)^2 + (11.9 - 12.0167)^2 + \ldots + (11.8 - 12.0167)^2 }{ 12 - 1 } \\
                &\approx 0.0452.\rubric{3}
        \end{align*}

        We berekenen het steekproefgemiddelde $\overline{x_B}$ voor de $m = 10$ levertijden van leverancier $B$ als volgt:
        \begin{align*}
            \overline{x_B} = \frac{b_1+b_2+\ldots+b_m}{m} = \frac{ 11.9 + 7.3 + \ldots + 11.4 }{ 10 } = 9.48.\rubric{1}
        \end{align*}
    

        We berekenen de steekproefvariantie $s_B^2$ als volgt:
        \begin{align*}
            s_B^2 &= \frac{ (b_1 - \overline{x_B})^2 + (b_2 - \overline{x_B})^2 + \ldots + (b_m - \overline{x_B})^2 }{ m - 1 } \\
                &= \frac{ (11.9 - 9.48)^2 + (7.3 - 9.48)^2 + \ldots + (11.4 - 9.48)^2 }{ 10 - 1 } \\\rubric{2}
                &\approx 10.7173.
        \end{align*}
    }

    \subquestion{10}{
        Voer een $F$-toets uit om te bepalen of de varianties in de levertijden gelijk kunnen zijn voor beide leveranciers.
        Geef hierbij expliciet de stappen aan van het algemene stappenplan voor hypothesetoetsen.
        Bepaal de toetsuitslag op basis van het kritieke gebied.
    }
    \solution{
        \subsubsection*{Stap 1: formuleer de nulhypothese $H_0$ en de alternatieve hypothese $H_1$.}
        Bij een $F$-toets toetsen we of de varianties $\sigma_A^2$ en $\sigma_B^2$ gelijk zouden kunnen zijn, oftewel
        \begin{align*}
            H_0: \quad \sigma_A^2 = \sigma_B^2 \\
            H_1: \quad \sigma_A^2 \neq \sigma_B^2 \rubric{1}
        \end{align*}

        \subsubsection*{Stap 2: bepaal het significantieniveau $\alpha$.}
        Er is in de vraag gegeven dat we mogen werken met een significantieniveau $\alpha=0.05$.
        
        \subsubsection*{Stap 3: verzamel data voor het berekenen van de toetsingsgrootheid.}
        Ook zijn er in de vraag al steekproefdata verzameld voor beide populaties. 
        Op basis van deze steekproefdata hebben we in vraag (a) bepaald dat:
        
        \begin{center}
            \begin{tabular}{C{.3\textwidth} C{.3\textwidth} C{.3\textwidth}}
                $\overline{x_A} \approx 12.0167$ & $s_A^2 \approx 0.0452$ & $n = 12$ \\
                $\overline{x_B} = 9.48$ & $s_B^2 \approx 10.7173$ & $m = 10$
            \end{tabular}
        \end{center}

        \subsubsection*{Stap 4: bereken de toetsingsgrootheid.}
        De theoretische toetsingsgrootheid (hoofdletter!) bij een $F$-toets is gelijk aan
        \[
            F = \frac{S_A^2}{S_B^2}, \rubric{1}
        \]
        en volgt onder $H_0$ een $F$-verdeling met $\dfn = n-1 = 11$ en $\dfd = m-1 = 9$ vrijheidsgraden. \rubric{1}

        De geobserveerde toetsingsgrootheid (kleine letter!) is gelijk aan
        \[
            f = \frac{s_A^2}{s_B^2} = \frac{0.0452}{10.7173} \approx 0.0042. \rubric{1}
        \]
        
        \subsubsection*{Stap 5: bepaal de toetsuitslag (met behulp van het kritieke gebied).}

        Een $F$-toets voor gelijke varianties is altijd tweezijdig, oftewel het kritieke gebied is van de vorm $(0, g_1)$ en $(g_2, \infty)$, waarbij de grenswaarden $g_1$ en $g_2$ als volgt kunnen worden berekend:
        \begin{enumerate}
            \item Bepaal de grenswaarde $g_1$ waarvoor geldt dat $P(F \leq g_1) = \alpha/2 = 0.025$.
            Voer hiervoor in de \enquote{numerieke solver} van de grafische rekenmachine in:
            \begin{align*}
                E_1 &= \fcdf(\lb=0, \ub=X, \df_1=11, \df_2=9) \\
                E_2 &= 0.025
            \end{align*}
            De \enquote{numerieke solver} geeft een waarde van $g_1 \approx 0.2787$.
            \item Bepaal de grenswaarde $g_2$ waarvoor geldt dat $P(F \geq g_2) = \alpha/2 = 0.025$.
            Voer hiervoor in de \enquote{numerieke solver} van de grafische rekenmachine in:
            \begin{align*}
                E_1 &= \fcdf(\lb=X, \ub=\GRinfinity, \df_1=11, \df_2=9) \\
                E_2 &= 0.025
            \end{align*}
            De \enquote{numerieke solver} geeft een waarde van $g_2 \approx 3.9121$.
        \end{enumerate}
        
        Het kritieke gebied is dus gelijk aan $(0, 0.2787)$ en $(3.9121, \infty)$. \rubric{4}

        De geobserveerde toetsingsgrootheid $f \approx 0.0042$ ligt in het kritieke gebied, dus we besluiten om $H_0$ te verwerpen.\rubric{1}

        \subsubsection*{Stap 6: vertaal de toetsuitslag terug naar de originele context van het probleem.}
        Er is voldoende reden om op basis van deze steekproeven de aanname van gelijke varianties (van de levertijden van de twee leveranciers) te verwerpen, oftewel we gaan er nu van uit dat $\sigma_X^2 \neq \sigma_Y^2$. \rubric{1}
        \begin{center}
            \resizebox{0.9\textwidth}{!}{
                \includegraphics{exam20260724-f-kritiek.png} 
            }
        \end{center}
    }

    \subquestion{10}{
        Bepaal met behulp van een geschikte $t$-toets of de gemiddelde levertijd significant verschilt tussen de twee leveranciers.
        Geef hierbij expliciet de stappen aan van het algemene stappenplan voor hypothesetoetsen.
        Bepaal hierbij de toetsuitslag met behulp van de $p$-waarde.
    }
    \solution{
        \subsubsection*{Stap 1: formuleer de nulhypothese $H_0$ en de alternatieve hypothese $H_1$.}
        We gaan toetsen of de gemiddelde levertijden $\mu_A$ en $\mu_B$ voor beide leveranciers significant van elkaar verschillen.
        De hypothesen zijn daarom als volgt:
        \begin{align*}
            H_0: &\quad \mu_A = \mu_B \text{ (geen significant verschil in de gemiddelde levertijd) } \\
            H_1: &\quad \mu_A \neq \mu_B \text{ (wel een significant verschil in de gemiddelde levertijd) } \rubric{2}
        \end{align*}

        \subsubsection*{Stap 2: bepaal het significantieniveau $\alpha$.}
        Er is opnieuw gegeven dat we mogen werken met een significantieniveau $\alpha=0.05$.
        
        \subsubsection*{Stap 3: verzamel data voor het berekenen van de toetsingsgrootheid}
        Ook zijn er in de vraag al steekproefdata verzameld voor beide populaties. 
        Op basis van deze steekproefdata hebben we in vraag (a) bepaald dat:
        
        \begin{center}
            \begin{tabular}{C{.3\textwidth} C{.3\textwidth} C{.3\textwidth}}
                $\overline{x_A} \approx 12.0167$ & $s_A^2 \approx 0.0452$ & $n = 12$ \\
                $\overline{x_B} = 9.48$ & $s_B^2 \approx 10.7173$ & $m = 10$
            \end{tabular}
        \end{center}

        \subsubsection*{Stap 4: bereken de toetsingsgrootheid.}
        Op basis van ons antwoord bij vraag (b) mogen we aannemen dat de varianties ongelijk zijn, oftewel $\sigma_X^2 \neq \sigma_Y^2$.
        De theoretische toetsingsgrootheid (hoofdletter!) bij een $t$-toets met twee steekproeven (met ongelijke varianties) is gelijk aan
        \[
            T = \frac{(\overline{X_A} - \overline{X_B}) - (\mu_A - \mu_B)}{\sqrt{\frac{S_A^2}{n} + \frac{S_B^2}{m}}}, \rubric{1}
        \]
        en volgt onder $H_0$ een $t$-verdeling met $\df=\min(n-1,m-1)=9$ vrijheidsgraden. \rubric{1}

        De geobserveerde toetsingsgrootheid (kleine letter!) is gelijk aan
        \begin{align*}
            t   &= \frac{(\overline{x_A}-\overline{x_B}) - (\mu_A - \mu_B)}{\sqrt{\frac{s_A^2}{n} + \frac{s_B^2}{m}}} \\
                &= \frac{(12.0167 - 9.48) - 0}{\sqrt{\frac{0.0452}{12} + \frac{10.7173}{10}}} \\
                &\approx 2.4460. \rubric{2}
        \end{align*}

        \subsubsection*{Stap 5: bepaal de toetsuitslag (met behulp van de $p$-waarde).}
        Omdat we tweezijdig toetsen, is de $p$-waarde gelijk aan het minimum van $P(T \leq t)$ en $P(T \geq t)$:
        Deze kansen kunnen we met de $t$-verdeling bepalen, namelijk:
        \begin{align*}
            P(T \leq t) &= \tcdf(\lb = -\GRinfinity, \ub=2.4460, \df=9) \approx 0.9815, \\ 
            P(T \geq t) &= \tcdf(\lb = 2.4460, \ub=\GRinfinity, \df=9) \approx 0.0185.
        \end{align*}
        De $p$-waarde is dus gelijk aan $p \approx 0.0185$. \rubric{2}

        Aangezien de $p$-waarde daarmee kleiner is dan $\alpha/2$ (tweezijdige toets!), besluiten we om de nulhypothese te verwerpen.\rubric{1}
        
        \subsubsection*{Stap 6: vertaal de toetsuitslag terug naar de originele context van het probleem.}
        Er is op basis van deze steekproeven voldoende reden om aan te nemen dat de gemiddelde levertijd significant verschilt 
        tussen de twee leveranciers. \rubric{1}    
        \begin{center}
            \resizebox{0.9\textwidth}{!}{
                \includegraphics{exam20260724-two-sample-t-test-pvalue.png} 
            }
        \end{center}
    }

    \subquestion{3}{
        Is er op basis van deze hypothesetoetsen een eenduidige voorkeur aan te wijzen voor één leverancier?
        Zo ja, aan welke leverancier zou je de voorkeur geven en waarom?
        Zo nee, geef voor beide leveranciers een reden die in het voordeel van deze leverancier spreekt.
    }
    \solution{
        Er is geen eenduidige voorkeur aan te wijzen voor de ene of de andere leverancier. \rubric{1}
        Uit de hypothesetoets bij vraag (c) is op te maken dat de gemiddelde levertijd van leverancier $B$ met grote betrouwbaarheid significant 
        korter is dan bij leverancier $A$ ($\mu_A < \mu_B$), wat in het voordeel van leverancier $B$ zou spreken.\rubric{1}
        Uit de $F$-toets voor gelijke varianties blijkt echter dat de variantie in de levertijden bij leverancier $B$ veel groter is ($\sigma_A^2 << \sigma_B^2$), waardoor
        leverancier $A$ als betrouwbaarder en voorspelbaarder gezien kan worden.\rubric{1}
    }
\end{question}